\DTMsavedate{mydate}{2022-06-30}
\maketitle{Digital Communication Engineering}{Electronic Engineering}{\DTMusedate{mydate}}

% ----------------------------------------------------- %
% COURSE INFO
\subsection*{Course Information}
\vspace{0ex}%
\noindent\parbox{0.5\textwidth}{%
    \noindent\begin{tabular}{@{}ll}
                 \textsf{Course:}          & Digital Communication Engineering 	 \\
                 \textsf{Course Code:}         & 	TEE 3221    \\
                 \textsf{Credit Hours:}    & 	10
    \end{tabular}}
\vspace{2ex}\hrule\vspace{2ex}

% ----------------------------------------------------- %
% INSTRUCTOR INFO
\subsection*{Lecturer Information}
\noindent\parbox{0.5\textwidth}{%
    \noindent\begin{tabular}{@{}ll}
                 \textsf{Name:}         &     Magripa Nleya         \\
%\textsf{Phone:}        &    \phone          \\
\textsf{Email:}        &    magripa.nleya@nust.ac.zw \\
\textsf{Qualifications:}        &  MSc in Telecommunication Engineering
    \end{tabular}}
\vspace{2ex}\hrule\vspace{2ex}


% ----------------------------------------------------- %
% LEARNING OBJECTIVES
\subsection*{Learning Objectives}
The course covers the GSM.
The course will make the students understand:
\begin{outline}
    \1 the principles of digital communications.

    \1 coding.

    \1 modulation and multiplexing methods.
\end{outline}
\vspace{2ex}\hrule\vspace{2ex}

% ----------------------------------------------------- %
% COURSE TOPICS
\subsection*{Topics}
\begin{outline}
    \1 Introduction to Digital Communication Systems (DSC)
    \begin{itemize}
        \item Spectral Density
        \item Autocorrelation
        \item Random Signals
    \end{itemize}
    \1 Formatting and Baseband Transmission
    \begin{itemize}
        \item Baseband Systems
        \item Formatting textual Data (character coding)
        \item Formatting Analog Information
        \item Sources of Corruption
        \item Pulse Code Modulation
        \item Uniform and Non-Uniform Quantization
        \item Baseband transmission
    \end{itemize}
    \1 Global System for Mobile Communications(GSM)
    \begin{itemize}
        \item The GSM subsystems
        \item The Network Subsystem
        \item The Base Station Subsystem (BSS)
        \item Mobility management and call control
        \item Call reselection and location area update
        \item The mobile-terminated call
        \item Handover scenarios
        \item The mobile device
        \item The SIM card
        \item The intelligent network subsystem and CAMEL
    \end{itemize}
\end{outline}
\vspace{2ex} \hrule\vspace{2ex}

% ----------------------------------------------------- %
% Students Assenssment
\subsection*{Assessment of Students}
\begin{outline}
    \1 Continuous assessment will consist of written assignments and written in-class tests.
    \1 The final assessment shall be an examination written at the end of the semester.
\end{outline}
\vspace{2ex}\hrule\vspace{2ex}

% ----------------------------------------------------- %
% Grade Evaluation
\subsection*{Course Evaluation and Grading Policies}
Grades are based on:
Continuous assessment (\qty{25}{\percent}),
Final Exam (\qty{75}{\percent})
\vspace{2ex}\hrule\vspace{2ex}
% ----------------------------------------------------- %
% EQUIPMENT
% https://sites.udel.edu/computing-purchases/personal-specs/
% \subsection*{Essential Equipment}

% \vspace{2ex}\hrule\vspace{2ex}

% Policies and Other information
\subsection*{Policies and Other Information}
% Key Text or some Books
\subsection*{Key Text}
\begin{itemize}
    \item Communications Engineering Principles by I. Otung
    \item Digital Communications: Fundamentals and Applications by B. Sklar
    \item From GSM to LTE by M. Sauter
    \item Lecture notes
\end{itemize}
\vspace{2ex}\hrule
\vspace{2ex}\hrule\vspace{2ex}